\section{Notation and formal definitions}
\label{app:notation}

\subsection*{Units and symbols}

\begin{longtable}[]{@{}ll@{}}
\toprule\noalign{}
Symbol & Meaning \\
\midrule\noalign{}
\endhead
$n$ & one official BOAMP notice \\
$e$ & one reconstructed procurement episode \\
$i$ & an \emph{anchor}: an awarded cohort episode \\
$j$ & a \emph{candidate}: a later episode exposed for comparison \\
$u_i$ & award date of anchor $i$ (time origin) \\
$v_j$ & first-publication date of candidate $j$ \\
$J_i$ & candidate set surviving blocking for anchor $i$ \\
$T_{ij}$ & text similarity between anchor $i$ and candidate $j$ \\
$\hat{j}_i$ & the accepted successor, or $\varnothing$ on abstention \\
$Y_i$ & event indicator, $1$ if a successor is accepted \\
$\tau_i$ & observed or censored time from $u_i$ \\
$S(t)$ & probability of no observable successor by $t$ \\
$h(t\mid X)$ & hazard of an observable successor given covariates $X$ \\
$Z_i$ & annotated technology class of a notice \\
$\hat{Z}_e$ & predicted technology class of a cohort episode \\
\bottomrule\noalign{}
\end{longtable}

\subsection*{Cohort rule}

An episode $e$ enters the study cohort if it is awarded, geolocated in the Grand Ouest, has a resolvable
award date, and
\[
\exists\, c \in \mathrm{CPV}(e) \quad\text{such that}\quad \left\lfloor c/10^{6} \right\rfloor \in \{32,35,48,72\}.
\]
This is an \emph{any-code} rule at episode level, not a rule about the main CPV. The stratifying variable
\texttt{digital\_segment} is the lowest-numbered qualifying division present, so each episode contributes to
exactly one curve.

\subsection*{Blocking rule (candidate generation)}

Candidate $j$ is exposed for anchor $i$ if buyer identifiers are compatible --- same buyer key or same
normalised buyer name, with two \emph{different} validated SIRENs disqualifying --- and
\[
u_i + 90\ \text{days} \;\le\; v_j \;\le\; u_i + 2{,}920\ \text{days}.
\]
Compatibility is deliberately asymmetric: two unknown SIRENs are \emph{not} treated as evidence of identity.

\subsection*{Buyer evidence ladder}

The strongest available evidence determines the buyer component: validated SIREN $=1.00$, exact normalised
name $=0.90$, fuzzy name match at ratio $\ge 0.82$ giving $0.70$, informative-token Jaccard $\ge 0.34$ giving
$0.60$, otherwise $0$. Municipal prefixes (\emph{commune de}, \emph{ville de}, \dots) are stripped for
blocking; intercommunal legal forms are preserved, so \emph{Rennes Métropole} does not collapse onto
\emph{Ville de Rennes}.

\subsection*{Decision rule (primary event definition)}

\[
\hat{j}_i = \arg\max_{j \in J_i} T_{ij},
\qquad
Y_i = \mathbf{1}\!\left(T_{i\hat{j}_i} \ge 0.70\right),
\]
where $T_{ij}$ is the maximum of a word-level TF--IDF cosine (1--2 grams) and a character-level TF--IDF
cosine (\texttt{char\_wb}, 3--5 grams). At most one successor is accepted per anchor; otherwise the method
abstains.

\subsection*{Missing evidence versus negative evidence}

Two components return \texttt{NaN} rather than $0$ when the evidence does not exist: CPV continuity when
either side carries no CPV code, and the time component when the anchor has no reliable declared duration.
The weighted score of the contrast arm renormalises over the components that do carry evidence,
\[
\text{score} = 100 \times \frac{\sum_{k \in \mathcal{O}} w_k \, c_k}{\sum_{k \in \mathcal{O}} w_k},
\qquad \mathcal{O} = \{k : c_k \text{ observed}\},
\]
so a pair lacking CPV evidence is judged on its remaining evidence rather than penalised as though the codes
disagreed.

\subsection*{Survival variables}

\[
\tau_i =
\begin{cases}
v_{\hat{j}_i} - u_i, & Y_i = 1,\\[2pt]
\text{2025-12-31} - u_i, & Y_i = 0 \quad \text{(administrative right censoring)}.
\end{cases}
\]

\subsection*{Procurement family (NLP grouping unit)}

Two annotated notices belong to the same family if they were placed in one reconstructed episode, or if
their object texts reach a character-level cosine similarity of at least $0.80$. Every family sits in exactly
one cross-validation fold.
